\documentclass{article}
\usepackage[UTF8]{ctex}
\usepackage{amsmath,amssymb,amsfonts}
\usepackage{enumitem}
\usepackage{geometry}
\geometry{a4paper, margin=1in}
% === 错题本排版优化（不参与题目标识） ===
\setlength{\parindent}{0pt}
\setlength{\parskip}{0.35em plus 0.10em minus 0.08em}
\setlength{\emergencystretch}{3em}
\setlist{itemsep=0.25em, topsep=0.40em, parsep=0pt}
\allowdisplaybreaks[2]
% === 错题本排版优化结束 ===
\begin{document}
\section*{南开大学2026年高等代数真题}
\section*{50．南开大学2026年高等代数试卷}
\begin{itemize}

\subsection{题目1}
  \item[1．] 求次数最小的多项式 \(f(x)\) ，使得
\[
(x+1)^{3}\left|f(x)-1,(x-1)^{2}\right| f(x)+1 .
\]

\subsection{题目2}
  \item[2．] 计算 \(n\) 阶行列式
\[
\left|\begin{array}{ccccc}
1^{2} & 2^{2} & 3^{2} & \cdots & n^{2} \\
2^{2} & 3^{2} & 4^{2} & \cdots & (n+1)^{2} \\
3^{2} & 4^{2} & 5^{2} & \cdots & (n+2)^{2} \\
\vdots & \vdots & \vdots & & \vdots \\
n^{2} & (n+1)^{2} & (n+2)^{2} & \cdots & (2 n-1)^{2}
\end{array}\right| .
\]
\end{itemize}

\subsection{题目3}
3．设

\[
A=\left(\begin{array}{lll}
1 & 1 & a \\
1 & a & 1 \\
a & 1 & 1
\end{array}\right), \beta=\left(\begin{array}{c}
1 \\
1 \\
-2
\end{array}\right) .
\]

如果方程组 \(A X=\beta\) 有无穷多个解。求 \(a\) 的值，并求正交矩阵 \(Q\) ，使得 \(Q^{\mathrm{T}} A Q\) 为对角矩阵。

\subsection{题目4}
4．设 \(\mathbb{R}\) 为实数域，\(A \in \mathbb{R}^{n \times n}, \beta, \gamma \in \mathbb{R}^{n \times 1}, a \in \mathbb{R}\) ，证明：

\[
\left|\begin{array}{cc}
A & \beta \\
\gamma^{\mathrm{T}} & a
\end{array}\right|=a|A|-\gamma^{\mathrm{T}} A^{*} \beta .
\]

其中 \(A^{*}\) 为 \(A\) 的伴随矩阵。

\begin{itemize}

\subsection{题目5}
\end{itemize}

\subsection{题目5}
原题边界“5．”：设 \(V_{1}, V_{2}, V_{3}\) 都是有限维线性空间 \(V\) 的子空间，证明：
\[
\operatorname{dim}\left(V_{1} \cap V_{2}\right)+\operatorname{dim}\left(\left(V_{1}+V_{2}\right) \cap V_{3}\right)=\operatorname{dim}\left(V_{2} \cap V_{3}\right)+\operatorname{dim}\left(\left(V_{2}+V_{3}\right) \cap V_{1}\right) .
\]

\subsection{题目6}

\subsection{题目6}
原题边界“6．”：设 \(A=\left(a_{i j}\right)_{n \times n}\) 为实正定矩阵，证明：\(|A| \leq a_{11} a_{22} \cdots a_{n n}\) ．

\subsection{题目7}
7．设 \(A\) 为 \(n\) 阶实对称矩阵，线性变换 \(\varphi: \mathbb{R}^{n \times n} \rightarrow \mathbb{R}^{n \times n}\) 定义为 \(\varphi(X)=A X A, \forall X \in \mathbb{R}^{n \times n}\) 。

\begin{itemize}
  \item[（1）] 若 \(\lambda_{1}, \lambda_{2}\) 是 \(A\) 的特征值，证明：\(\lambda_{1} \lambda_{2}\) 是 \(\varphi\) 的特征值．
  \item[（2）] 证明：存在 \(\mathbb{R}^{n \times n}\) 的一组基，使得 \(\varphi\) 在该基下的矩阵为对角矩阵。
\end{itemize}

\subsection{题目8}
8．已知 \(n\) 阶复方阵 \(A\) 与 \(B\) 的秩相等，且 \(A^{2} B=A\) ，证明：\(B^{2} A=B\) 。
\end{document}
